%! Vectors — Unit 8, Lesson 8.6 — Guided Notes
\documentclass[10pt]{article}
\usepackage{precalculus-article}
\usepackage{precalculus-boxes}
\newcommand{\TallMath}[1]{$\displaystyle #1\rule[-1.4em]{0pt}{3.2em}$}

\begin{document}

\pageheader{Unit 8, Lesson 8.6}{Guided Notes}
\namedateperiod

\vspace{0.1in}

\begin{objectivebox}
By the end of this lesson, I will be able to:
\begin{itemize}[leftmargin=1.5em, itemsep=3pt]
  \item \writeline
  \item \writeline
  \item \writeline
  \item \writeline
\end{itemize}
\end{objectivebox}

\vspace{0.08in}

\begin{vocabbox}
\termblanklong{vector}
\termblanklong{tail / head}
\termblanklong{magnitude}
\termblanklong{scalar multiplication}
\termblanklong{dot product}
\termblanklong{unit vector}
\end{vocabbox}

\vspace{0.08in}

\begin{hookbox}
Two drones receive identical displacement instructions. Drone A leaves $(1,2)$ and
arrives at $(5,6)$. Drone B leaves $(3,0)$ and arrives at $(7,4)$.
How far east and how far north does each drone travel? Why do they receive the same instructions?

\writeline
\end{hookbox}

\vspace{0.08in}

\begin{notesbox}{Section 1 --- Vector Components and Magnitude}
\small
\textbf{Definition.} A \textbf{vector} from tail $P_1=(x_1,y_1)$ to head $P_2=(x_2,y_2)$ has
component form $\vec{v} = \langle a, b\rangle$ where:

\vspace{0.06in}
\hspace{1em} $a = x_2 - x_1 = $ \blank{3cm} \qquad $b = y_2 - y_1 = $ \blank{3cm}

\vspace{0.08in}
\textbf{Magnitude:} $\|\vec{v}\| = \sqrt{a^2 + b^2} = $ \blank{5cm}

\vspace{0.12in}
\textbf{Direction from trig.} If $\|\vec{v}\| = r$ and direction angle is $\theta$:

\vspace{0.06in}
\hspace{1em} $a = r\cos\theta = $ \blank{3.5cm} \qquad $b = r\sin\theta = $ \blank{3.5cm}

\vspace{0.12in}
\textbf{Example.} Tail $P_1=(1,3)$, head $P_2=(5,6)$.

\vspace{0.06in}
\hspace{1em} $a = $ \blank{2cm} \quad $b = $ \blank{2cm} \quad
$\vec{v} = \langle$ \blank{1.5cm}$,$ \blank{1.5cm}$\rangle$ \quad
$\|\vec{v}\| = $ \blank{2.5cm}
\end{notesbox}

\vspace{0.08in}

\begin{notesbox}{Section 2 --- Scalar Multiplication and Vector Addition}
\small
\textbf{Scalar multiplication.} $k\langle a,b\rangle = \langle$ \blank{3cm}$\rangle$

\vspace{0.06in}
New magnitude: $|k| \cdot \|\vec{v}\|$. \quad
Direction: \blank{3cm} when $k>0$; \blank{3cm} when $k<0$.

\vspace{0.12in}
\textbf{Vector addition.} $\langle a_1,b_1\rangle + \langle a_2,b_2\rangle
= \langle$ \blank{4cm}$\rangle$

\vspace{0.06in}
Geometric interpretation: place tail of second vector at \blank{3cm} of first;
resultant runs from \blank{3cm} to \blank{3cm}.

\vspace{0.12in}
\textbf{Example.} $\vec{u}=\langle 3,1\rangle$, $\vec{w}=\langle -1,4\rangle$.

\vspace{0.06in}
\renewcommand{\arraystretch}{1.6}
\begin{tabular}{lll}
$\vec{u}+\vec{w} = \langle$\blank{1.5cm}$,$\blank{1.5cm}$\rangle$ &\quad&
$2\vec{u} = \langle$\blank{1.5cm}$,$\blank{1.5cm}$\rangle$ \\
\end{tabular}
\end{notesbox}

\vspace{0.08in}

\begin{notesbox}{Section 3 --- Dot Product and Unit Vectors}
\small
\textbf{Dot product.}
$\langle a_1,b_1\rangle \cdot \langle a_2,b_2\rangle
= a_1 a_2 + b_1 b_2 = $ \blank{4cm}
\quad (result is a \blank{2.5cm}, not a vector)

\vspace{0.12in}
\textbf{Unit vector.} A vector of magnitude \blank{0.8cm}.

\vspace{0.06in}
$\hat{u} = \dfrac{\vec{v}}{\|\vec{v}\|}
= \left\langle \dfrac{a}{\|\vec{v}\|},\, \dfrac{b}{\|\vec{v}\|} \right\rangle$

\vspace{0.1in}
\textbf{Standard basis vectors:}
$\hat{\imath} = \langle$ \blank{2cm}$\rangle$ \quad
$\hat{\jmath} = \langle$ \blank{2cm}$\rangle$

Any vector: $\langle a,b\rangle =$ \blank{1.5cm}$\hat{\imath}\, +\, $\blank{1.5cm}$\hat{\jmath}$

\vspace{0.12in}
\textbf{Example.} $\vec{v}=\langle 3,4\rangle$, $\|\vec{v}\| = $ \blank{1.5cm}

\vspace{0.06in}
$\hat{u} = \left\langle \dfrac{3}{\blank{0.8cm}},\, \dfrac{4}{\blank{0.8cm}} \right\rangle
= \langle$\blank{2cm}$,$\blank{2cm}$\rangle$ \quad
Verify: $\|\hat{u}\| = $ \blank{2.5cm}
\end{notesbox}

\vspace{0.08in}

\begin{notesbox}{Section 4 --- Angle Between Vectors \& Perpendicularity}
\small
\textbf{Angle formula.}
$\vec{v}_1 \cdot \vec{v}_2 = \|\vec{v}_1\|\,\|\vec{v}_2\| \cos\theta$

\vspace{0.08in}
Solving for $\theta$:
\quad $\theta = \cos^{-1}\!\left( \dfrac{\vec{v}_1 \cdot \vec{v}_2}{\|\vec{v}_1\|\,\|\vec{v}_2\|} \right)
= $ \blank{5cm}

\vspace{0.12in}
\textbf{Perpendicular vectors.} Two nonzero vectors are perpendicular if and only if their
dot product $= $ \blank{1.5cm} \quad (because $\cos 90^\circ = $ \blank{1cm}).

\vspace{0.12in}
\textbf{Example.} $\vec{u}=\langle 1,2\rangle$, $\vec{w}=\langle 4,-2\rangle$.

\vspace{0.06in}
\hspace{1em} $\vec{u}\cdot\vec{w} = (1)(4) + (2)(-2) = $ \blank{2cm}

\hspace{1em} Are $\vec{u}$ and $\vec{w}$ perpendicular? \blank{3cm}

\vspace{0.1in}
\textbf{Law of Cosines application.} If $\vec{r} = \vec{u}+\vec{w}$, the triangle with sides
$\|\vec{u}\|$, $\|\vec{w}\|$, $\|\vec{r}\|$ satisfies:
$\|\vec{r}\|^2 = \|\vec{u}\|^2 + \|\vec{w}\|^2 - 2\|\vec{u}\|\,\|\vec{w}\|\cos\theta$

where $\theta$ is the angle \blank{5cm}.
\end{notesbox}

\end{document}
